\input{preamble}

\begin{document}

\title{Exact Marginal--Regret Frontiers\\for Signed-Prefix Compression}

\author{\IEEEauthorblockN{Vu Phan An Nguyen}
\IEEEauthorblockA{\textit{Rider University} \\
Lawrenceville, NJ, USA}}

\maketitle

\begin{abstract}
One random removal order may have to serve every compression budget, while each item must
have nearly uniform probability of removal at every prefix. We determine the price of this
constraint for a signed-weight family, using regret against the best subset at each budget.
For every size $m\ge3$ and every marginal-band width, we obtain the exact minimax expected regret.
With even size $m$, it is
$\max\{2/[(m+4)(m-1)],\,4(1-\eta)/[m(m+4)]\}$.
With odd size it has two affine pieces before reaching the unconstrained value at $\eta=1$.
Exact uniformity costs less than a factor two in even size but a factor proportional to $m$
in odd size. Explicit distributions attain the bounds simultaneously at all budgets; a
composition-flow formulation preserves both the objective and every item marginal.
The timing of the budget changes the answer: choosing it after the order is drawn
multiplies the free value by $m-1$, but the marginal constraint then costs at most $1+1/m$.
On the family and on $(a,-a,b,-b)$, knowing the budget before acting helps only once the band
is loose. On the family it gains
nothing up to the even breakpoint and the odd crossing point; on four items it gains nothing up
to $\eta=1/3$, and one exchange after the budget is revealed recovers the whole gain.
For arbitrary weight perturbations we prove a Lipschitz bound with recomputed baselines and
normalisation. In odd size, the infimum price ratio over an $\ell_1$ ball of radius
$0<\varepsilon\le1/6$ is $\Theta(\min\{m,1/\varepsilon\})$, with absolute constants.
This is a sharp scale for the erosion of a uniform guarantee, not a claim that every
perturbation destroys the parity effect. The core frontier, attainability, flow and stability
results have Lean~4 proofs, and so does the robust-ball statement.
For every multiset with $d$ weight types the value is that of a flow program with
$\prod_\tau(n_\tau+1)$ states per layer. The band is free exactly from an attained threshold
$\eta_{\rm free}$ computed by that program, and the additive price, unlike the ratio, is stable
with no lower bound on the free value.
Stored optimisation certificates and perturbation results provide finite-instance checks.
\end{abstract}

\begin{IEEEkeywords}
ranking under constraints, prefix regret, linear programming duality, formal verification
\end{IEEEkeywords}

\section{Introduction}
\label{sec:intro}

\subsection{The question}

A served compression must often work at a size that is not known in advance. One way to
serve it is as a \emph{removal order}. The system returns an order of the items, and a
consumer who must shed $k$ of them drops the first $k$ and keeps the other $m-k$. Fairness or
exposure requirements then constrain the distribution of the order: each item should be
dropped about equally often at every budget. We ask what that constraint costs when the order
must be good at every budget simultaneously.

Fix weights $g=(g_{1},\dots,g_{m})$ and a normaliser $D>0$. For an order $\pi$ and a budget
$k$, the best achievable loss and the regret are
\begin{equation}
b_{k}(g)=\min_{|S|=k}\Bigl|\sum_{i\in S}g_{i}\Bigr|,\qquad
r_{k}(\pi;g)=\frac{\bigl|\sum_{\ell\le k}g_{\pi_{\ell}}\bigr|-b_{k}(g)}{D}.
\label{eq:regret}
\end{equation}
Here $\{\pi_{1},\dots,\pi_{k}\}$ is the \emph{dropped} set. In the physical model of the
compression, keeping the complement of a set $S$ costs $|\sum_{i\in S}g_{i}|/2$. So with
$D=\|g\|_{1}$, the regret~\eqref{eq:regret} is the loss of the compression that keeps the
last $m-k$ items, above the best compression of that size, normalised by $\|g\|_{1}/2$.

\leanbox{\lean{physical\_radius}, \lean{physical\_minimum}, \lean{physical\_normalized\_regret}
in \lean{LeanPaperCStabilityEndgameV1}}

For a distribution $P$ over orders, the constrained optimum is
\begin{equation}
V_{g}(\eta)=\inf_{P}\ \max_{0<k<m}\ \mathbb{E}_{\pi\sim P}\,r_{k}(\pi;g)
\label{eq:value}
\end{equation}
subject to the marginal band
\begin{equation}
\Bigl|\Pr\nolimits_{P}\bigl(i\in\{\pi_{1},\dots,\pi_{k}\}\bigr)-\frac{k}{m}\Bigr|\le\frac{\eta}{m}
\quad(\text{all }i,\ k\le m).
\label{eq:band}
\end{equation}
The band reads the same whether it is stated for dropping or for keeping. An item is kept at
size $m-k$ with probability within $\eta/m$ of $\frac{m-k}{m}$ exactly
when~\eqref{eq:band} holds. The budget in~\eqref{eq:value} is fixed before the draw, so the
objective is $\max_{k}\mathbb{E}\,r_{k}$. Section~\ref{sec:expost} treats the objective
$\mathbb{E}\max_{k}r_{k}$, where the budget is chosen after the draw. The two objectives are
never interchanged.

\subsection{Results}

We work on the family $g_{m}$ of~\eqref{eq:family} below. We write
$V_{m}(\eta)=V_{g_{m}}(\eta)$ and $V^{\mathrm{free}}_{m}=\frac{2}{(m+4)(m-1)}$. The latter is
the unconstrained optimum, because the band is vacuous for $\eta\ge m$.

\begin{enumerate}
\item[1)] \textbf{The whole frontier} (Theorems~\ref{thm:zero}, \ref{thm:even},
\ref{thm:odd}). $V_{m}(\eta)$ in closed form for every $m\ge3$ and every $\eta\ge0$. At
$\eta=0$ the price $V_{m}(0)/V^{\mathrm{free}}_{m}$ is $\frac{2(m-1)}{m}<2$ in even size and
$\frac{(m+3)(m-1)}{2m}$ in odd size. The premium over one vanishes exactly at $\eta_{*}=\frac{m-2}{2(m-1)}$
(even) and at $\eta_{*}=1$ (odd).
\item[2)] \textbf{A polynomial-size program with the same value} (Theorem~\ref{thm:flow}).
\item[3)] \textbf{Robustness, and where it ends} (Theorems~\ref{thm:robust},
\ref{thm:collapse}, Corollary~\ref{cor:robustball}). A Lipschitz envelope for arbitrary
weights, and a two-sided law for the smallest price ratio in an $\ell_1$ ball:
$\Theta(\min\{m-1,1/\varepsilon\})$ for $0<\varepsilon\le1/6$.
Radius $1/m$ marks when a linear-in-$m$ price ceases to be guaranteed uniformly;
some much larger perturbations preserve it exactly.
\item[4)] \textbf{Any number of types} (Theorems~\ref{thm:typedflow}, \ref{thm:freeband},
Proposition~\ref{prop:gap}). For every multiset with $d$ weight types the flow program has the
same value over $\prod_\tau(n_\tau+1)$ states per layer. The band is free exactly from a threshold
$\eta_{\rm free}$, which is attained and read off that program. The additive price
$V_g(\eta)-V^{\rm free}_g$ is $6/\nu$-Lipschitz in $\ell_1$ on multisets with $\|g\|_1\ge\nu>0$, with
no lower bound on the free value.
On $(a,-a,b,-b)$ the normalised frontier $V_g/V_g^{\rm free}$ is the same at every ratio. A floating-point map of larger
two-magnitude multisets shows the ratio moving the breakpoint.
\item[5)] \textbf{Timing} (Proposition~\ref{prop:expost}, Theorem~\ref{thm:expostband}).
Choosing the budget after the draw multiplies the free value by $m-1$. Under the band the
ex-post value is known exactly, and there the band costs at most a factor $1+\frac1m$.
\item[6)] \textbf{Knowing the budget, and one exchange} (Proposition~\ref{prop:sandwich},
Theorems~\ref{thm:zerovalue}, \ref{thm:fourrecourse}). A policy that learns the budget first
is never worse, and $\lfloor m/2\rfloor$ exchanges after the reveal reach it. On $g_m$ it gains
nothing up to the even breakpoint and the odd crossing point. On $(a,-a,b,-b)$ it gains nothing
up to $\eta=\frac13$, and one exchange recovers the whole gain.
\end{enumerate}

\subsection{Prior work, and what does not follow from it}

Randomised rankings with constrained position or exposure marginals are a studied object.
Singh and Joachims~\cite{singh2018} introduce exposure constraints and optimise a linear
utility over the Birkhoff polytope. Celis, Straszak and Vishnoi~\cite{celis2018} bound group
representation in every prefix. Singh, Kempe and Joachims~\cite{singh2021} constrain, for every
item and every $k$, the probability that the item is ranked in the top $k$, and contrast ex-ante
and ex-post guarantees, as Freeman, Shah and Vaish~\cite{freeman2020} do in allocation. We
take that framing as given. We do not claim the flow technique either: reducing a permutation
polytope to a compositions DAG is standard~\cite{schrijver2003}. One ordering that serves every
budget is the object of incremental maximization, with competitive ratios for monotone
objectives~\cite{disser2023,bienkowski2026}. Choosing a dependence structure under fixed
marginals so that sums over unknown subsets cancel is studied through joint
mixes~\cite{koike2022}.

What is new is the law that the reduction lets us prove, for an objective these works do not
treat. The signed-prefix regret is an absolute sum net of the per-budget optimum, taken at
every budget. It is not linear in the order, and the parity of $k$ enters through $b_{k}$. It
depends on the joint law of each prefix, not only on the position marginals. With weights
$(+1,+1,-1,-1)$ on items $1,2,3,4$, the four cyclic shifts of the order $(1,3,2,4)$, with signs
$(+,-,+,-)$, and the four cyclic shifts of $(1,2,3,4)$, with signs $(+,+,-,-)$, each put every item
at every position with probability $\frac14$. Their expected $|\text{prefix}_{2}|$ are $0$ and
$1$. So the marginal rank matrix, which suffices for a linear utility, does not
suffice here. Against incremental maximization, the values here are exact minimax values under
marginal bands with signed cancellation and a per-budget optimum, not competitive ratios. The
joint-mix theory does not restrict realisations to permutations of one multiset, nor subsets to
nested prefixes, and it does not give the frontiers below. Recourse, the number of past
decisions an online algorithm may revise, is studied for discrepancy~\cite{gupta2022}. Here the
revision happens once, after the budget is revealed, and its value is exact.

\begin{table}[h]
\caption{Closest research threads and the specific quantitative results compared here.}
\label{tab:prior}
\centering
\footnotesize
\begin{tabular}{@{}p{0.30\columnwidth}p{0.62\columnwidth}@{}}
\toprule
Nearest source & Specific result compared here \\
\midrule
Exposure-constrained rankings~\cite{singh2018} & The exact signed-prefix frontier at every
size and band, with a $\Theta(m)$ parity split (Thms.~\ref{thm:zero}--\ref{thm:odd}) \\
Prefix-fair ranking~\cite{celis2018} & The exact breakpoints and the crossing point
$\eta_{t}$ of the odd frontier (Thm.~\ref{thm:odd}) \\
Ex-ante vs.\ ex-post~\cite{singh2021,freeman2020} & The exact banded ex-post value, and the
fact that the band and timing prices do not multiply (Thm.~\ref{thm:expostband}) \\
Flow reductions~\cite{schrijver2003} & Value equality with every per-item band carried
through relabelling, for every typed multiset over $\prod_\tau(n_\tau+1)$ states, and the free
threshold read off it (Thms.~\ref{thm:flow}, \ref{thm:typedflow}, \ref{thm:freeband}) \\
Incremental maximization~\cite{disser2023,bienkowski2026} & Exact values, not ratios, under
marginal bands and signed cancellation (Thms.~\ref{thm:zero}--\ref{thm:odd}) \\
Joint mixes~\cite{koike2022} & The prefix-nested, permutation-restricted frontiers
(Thms.~\ref{thm:zero}--\ref{thm:odd}) \\
Recourse~\cite{gupta2022} & The exact value of knowing the budget, and of one exchange after it
is revealed, under marginal bands (Thms.~\ref{thm:zerovalue}, \ref{thm:fourrecourse}) \\
\bottomrule
\end{tabular}
\end{table}

\subsection{What is proved, and how}

Three evidence classes appear, and they are never interchanged.
\begin{itemize}
\item \PROVED: a theorem stated at the paper's objects and checked in
Lean~4~\cite{demoura2021} against Mathlib~\cite{mathlib2020}.
\item \CERT: an exact rational certificate verified by a script.
\item \NUM: a floating-point solve or search.
\end{itemize}
Each formal theorem carries a line naming its Lean declaration. The robust-ball corollary
is proved in this text and checked as \lean{robust\_ball}. The proof architecture and supporting
declarations are in the supplement. Audited modules use no \texttt{sorry} or
\texttt{native\_decide}, and their axioms lie within \{\texttt{propext},
\texttt{Classical.choice}, \texttt{Quot.sound}\}. A compiled theorem certifies its statement
under its hypotheses. Every hypothesis of a headline is either a model assumption displayed
in it or discharged.

\section{Setting}
\label{sec:setting}

The family studied throughout is
\begin{equation}
g_{m}=\bigl(-3,\underbrace{-1,\dots,-1}_{p},\underbrace{1,\dots,1}_{q},3\bigr),
\qquad D=m+4,
\label{eq:family}
\end{equation}
with $p=\lceil\frac{m-2}{2}\rceil$ and $q=\lfloor\frac{m-2}{2}\rfloor$. The total is $0$ for
even $m$ and $-1$ for odd $m$, and $\|g_{m}\|_{1}=m+4=D$. Every weight is an odd integer, so
a sum of $k$ weights has the parity of $k$ and has absolute value at least $k\bmod2$. For
$g_{m}$ the per-budget optimum is exactly that,
\begin{equation}
b_{k}(g_{m})=k\bmod 2\qquad(0<k<m),
\label{eq:baseline}
\end{equation}
attained by an explicit subset at each budget. So on $g_{m}$ the physical regret, with the
baseline computed from its definition, is the normalised parity excess.

\leanbox{\lean{familyBaselineIdentification\_holds} in \lean{LeanPaperCBaselineAttainV1};
\lean{trueValue\_family} in \lean{LeanPaperCStabilityFamilyV1}}

The parity split lives on this structure. There are two heavy items of weight $\pm3$, and
light items of weight $\pm1$ whose counts differ by one in odd size. A $+3$ step between two
budgets of zero excess is impossible, because a zero-excess prefix sum has absolute value at
most $1$. That observation drives every odd lower bound below.

\section{The frontier}
\label{sec:frontier}

\begin{theorem}[Zero band]
\label{thm:zero}
For every even $m\ge4$ and every odd $m\ge3$,
\[
V_{m}(0)=
\begin{cases}
\dfrac{4}{m(m+4)}, & m\text{ even},\\[8pt]
\dfrac{m+3}{m(m+4)}, & m\text{ odd},
\end{cases}
\]
attained by an explicit law on all $m!$ orders of $g_{m}$. Hence
$V_{m}(0)/V^{\mathrm{free}}_{m}$ is $\frac{2(m-1)}{m}$ for even $m$ and
$\frac{(m+3)(m-1)}{2m}$ for odd $m$.
\end{theorem}

\leanbox{\lean{even\_value\_zero}, \lean{odd\_value\_zero}, \lean{odd\_uniform\_price\_ratio}
in \lean{LeanPaperCEndgameV1}; \lean{even\_allPerm\_witness} in \lean{LeanPaperCEvenWitnessJoinV1};
\lean{odd\_allPerm\_witness} in \lean{LeanPaperCOddWitnessJoinV1}}

\begin{theorem}[The even frontier]
\label{thm:even}
For every even $m\ge4$ and every $\eta\ge0$,
\[
V_{m}(\eta)=\max\Bigl\{\frac{2}{(m+4)(m-1)},\ \frac{4(1-\eta)}{m(m+4)}\Bigr\},
\]
so the price vanishes exactly at $\eta_{*}=\frac{m-2}{2(m-1)}$.
\end{theorem}

\leanbox{\lean{even\_frontier\_all} in \lean{LeanPaperCEndpointLawV1}}

\begin{theorem}[The odd frontier]
\label{thm:odd}
For every odd $m\ge3$ and every $\eta\ge0$,
\begin{multline*}
V_{m}(\eta)=\max\Bigl\{\frac{m+3-(m+1)\eta}{m(m+4)},\\
\frac{2}{(m+4)(m-1)}+\frac{(1-\eta)_{+}}{m(m+4)}\Bigr\}.
\end{multline*}
The two lines cross at $\eta_{t}=1-\frac{2}{m(m-1)}$, where
$V_{m}(\eta_{t})=\frac{2(m^{2}+1)}{m^{2}(m-1)(m+4)}$. In particular
$V_{m}(\eta)=V^{\mathrm{free}}_{m}$ exactly when $\eta\ge1$, and for any fixed $\eta<1$ the
odd price is $\Omega(m)$.
\end{theorem}

\leanbox{\lean{odd\_frontier}, \lean{odd\_kink\_value} in \lean{LeanPaperCOddFrontierV1};
\lean{odd\_breakpoint} in \lean{LeanPaperCEndpointLawV1}}

\emph{Why the frontier has this shape.} Every lower side comes from one budget or from an
average of budgets. Every upper side is one explicit law.

\emph{Zero band.} For even $m$ the weights sum to zero, so the prefix of length $m-1$ is minus
the last item. Then $r_{1}+r_{m-1}$ equals $\frac{2}{D}$ times the number of heavy items at
the two ends, and each end is heavy with probability $\frac{2}{m}$. For odd $m$ the prefix
of length $m-1$ is $-1$ minus the last item, and $m-1$ is even, so an item of weight $w$ in
last position costs $\frac{|1+w|}{D}$. Averaging over uniform last positions gives
$\frac{m+3}{mD}$. The upper laws are the following, each relabelled uniformly within weight
classes. For even $m$, the insertion mixture with one flat order. For odd $m$, the uniform
law on the cyclic shifts of $(+3,-3,-1,+1,\dots,-1)$.

\emph{The first odd line.} Only budget $m-1$ is used. By the band at $k=m-1$, the
last-position law satisfies $|\pi_{i}-\frac1m|\le\frac{\eta}{m}$. Hence for any centre $a$,
$\sum_{i}\pi_{i}c_{i}\ge\frac1m\sum_{i}c_{i}-\frac{\eta}{m}\sum_{i}|c_{i}-a|$. With
$c_{i}=|1+g_{i}|$ and $a=2$ the two sums are $m+3$ and $m+1$.

\emph{The second odd line.} Every order pays total excess at least $2+2[\pi_{m}>0]$. The
$+3$ step forces excess $2$ on the two budgets around it, and a positive last item forces two
more at budget $m-1$. The band puts each of the $\frac{m-1}{2}$ positive items last with
probability at least $\frac{1-\eta}{m}$. So the average over budgets, and hence the maximum,
is at least $V^{\mathrm{free}}_{m}+\frac{1-\eta}{mD}$.

\emph{The upper side.} $V_{m}$ is convex in $\eta$, so three laws suffice:
\begin{itemize}
\item the zero-band law at $\eta=0$;
\item the insertion mixture at $\eta=1$, whose value is $V^{\mathrm{free}}_{m}$;
\item a law at the crossing point $\eta_{t}$, with value exactly
$\frac{2(m^{2}+1)}{m^{2}(m-1)D}$.
\end{itemize}
The last law is new. For $m\ge7$ it is a mixture of explicit orders, with masses in units of
$\frac{1}{m^{2}(m-1)}$:
\begin{itemize}
\item mass $m^{2}+1$ on each insertion slot $0,\dots,m-3$ ($m^{2}-m+4$ at slot $0$), split
equally between two versions. Version~A is the insertion order. Version~B is its negation
with the last item set to $-1$;
\item $m(m-1)$ on version~A at slot $m-2$;
\item $2$ on the order that puts the light items of slot $m-2$ first and ends with
$(-3,+3)$;
\item $m-3$ on version~A at slot $0$ with its last two light items swapped, which ends with
$(-1,+1)$.
\end{itemize}
Every budget then carries expected excess exactly $\frac{2(m^{2}+1)}{m^{2}(m-1)}$, and the
class counts of every prefix satisfy the band at $\eta_{t}$. At $m=3$ and $m=5$ an explicit
finite law takes its place.

\section{The flow program has the same value}
\label{sec:flow}

Enumerating orders is hopeless past $m=11$. After symmetrisation, a prefix is characterised,
both for the objective and for the band, by its composition
$v=(n_{-3},n_{-1},n_{+1},n_{+3})\le(1,p,q,1)$. There are $4(p+1)(q+1)=O(m^{2})$ of them. Let
$\mathcal{F}_{m}(\eta)$ be the unit flows on the layered DAG of compositions whose edges add
one item, with the band imposed on the class occupancies.

\begin{theorem}[Flow equivalence]
\label{thm:flow}
Fix $m\ge2$, $\eta$, and a functional $\Phi$ of the budget vector $(\mathbb{E}r_{k})_{k}$. The
infimum of $\Phi$ over band-feasible laws on all $m!$ orders of $g_{m}$ equals its infimum
over $\mathcal{F}_{m}(\eta)$. In particular $V_{m}(\eta)$ is the value of the flow program,
and Theorems~\ref{thm:even} and~\ref{thm:odd} hold for it.
\end{theorem}

\leanbox{\lean{actual\_banded\_infima} in \lean{LeanPaperCFlowBandJoinV1};
\lean{odd\_frontier} in \lean{LeanPaperCOddFrontierV1}; \lean{support\_box} in
\lean{LeanPaperCFreeBandV1}; \lean{value\_eq\_flow}, \lean{even\_frontier\_flow}, \lean{odd\_breakpoint\_flow} in
\lean{LeanPaperCFlowObjectiveV1}}

Three maps make it exact.
\begin{itemize}
\item A law induces the occupancy of each composition. Conservation is the consistency of the
marginals of one joint law.
\item Averaging over relabellings inside a weight class keeps every $r_{k}$. It replaces each
item's marginal by its class average.
\item A unit flow carries a Markov law on whole trajectories. A trajectory of positive mass is
a sequence of one-item steps, hence the layer map of an actual order. Relabelling within
classes then recovers the per-item band.
\end{itemize}
The two programs reach exactly the same set of budget vectors, and that set is non-empty for
every $\eta\ge0$. The statement concerns the value of the four-type program on $g_{m}$. It
gives no algorithm for perturbed weights.

\begin{table}[t]
\caption{Flow-program sizes against the order formulation. \NUM{}}
\label{tab:agree}
\centering
\footnotesize
\begin{tabular}{@{}rrrl@{}}
\toprule
$m$ & flow states & flow edges & orders \\
\midrule
$11$ & $120$ & $316$ & $13{,}860$ type sequences \\
$15$ & $224$ & $612$ & intractable \\
$21$ & $440$ & $1236$ & intractable \\
$31$ & $960$ & $2756$ & intractable \\
\bottomrule
\end{tabular}
\end{table}

\section{Any number of weight types}
\label{sec:typed}

Theorem~\ref{thm:flow} uses only one fact about $g_m$: its items fall into weight classes. Let
$g=\mathrm{wt}\circ\mathrm{ty}$ for a type map $\mathrm{ty}:[m]\to[d]$ and type weights
$\mathrm{wt}\in\mathbb R^d$, and let $n_\tau$ be the multiplicity of type $\tau$. A count state
$v\in\prod_\tau[0,n_\tau]$ records how many items of each type the prefix has dropped. Its cost
at budget $k$ is $(|\sum_\tau\mathrm{wt}_\tau v_\tau|-b_k(g))/\|g\|_1$.

\begin{theorem}[The flow program for every typed multiset]
\label{thm:typedflow}
For every type map, weights and $\eta$:
\begin{itemize}
\item on a ranking, the count cost of its prefix composition is its normalised regret;
\item the budget vectors of band-feasible laws over all $m!$ rankings are exactly those of the
count-feasible banded flows, so every functional of the budget vector has the same infimum;
\item a count-feasible flow lives on the box $\prod_\tau[0,n_\tau]$, which has
$\prod_\tau(n_\tau+1)\le(m+1)^d$ states.
\end{itemize}
For a fixed number of types the program therefore has $O(m^d)$ states per layer. This is a
polynomial model for each fixed $d$, not a polynomial algorithm for arbitrary weights.
\end{theorem}
\leanbox{\lean{typedCost\_eq}, \lean{typed\_flow\_equivalence}, \lean{typed\_banded\_infima},
\lean{card\_state}, \lean{family\_is\_typed} in \lean{LeanPaperCTypedFlowV1};
\lean{support\_box}, \lean{card\_box} in \lean{LeanPaperCFreeBandV1}}

Flow reductions of permutation problems are not new~\cite{schrijver2003}. The content here is
that a nonlinear prefix objective and every per-item band survive the reduction together.

In this section and in Section~\ref{sec:aware} the normaliser is $D=\|g\|_1$.

\paragraph{When the band is free} The deviation of a law is
$\mathrm{dev}(P)=m\max_{k\le m,\,i}|P(i\text{ in the first }k)-k/m|$, so a law is band-feasible
at $\eta$ exactly when $\mathrm{dev}(P)\le\eta$. The band is vacuous at $\eta=m$, so
$V^{\rm free}_g=V_g(m)$. Let
\[
\eta_{\rm free}(g)=\min\{\mathrm{dev}(P):P\ \text{attains}\ V^{\rm free}_g\}.
\]

\begin{theorem}[The free threshold]
\label{thm:freeband}
Let $\|g\|_1>0$. The value and the threshold are attained. For every $\eta\ge0$,
$V_g(\eta)=V^{\rm free}_g$ if and only if $\eta\ge\eta_{\rm free}(g)$. For a typed multiset
the same holds with $V_g$ replaced by the value of the flow program of
Theorem~\ref{thm:typedflow}. So $\eta_{\rm free}$ is the least $\eta$ at which that program
reaches its value at $\eta=m$.
\end{theorem}
\leanbox{\lean{feasible\_iff\_dev}, \lean{feasible\_m\_eq}, \lean{attain\_value},
\lean{attain\_eta}, \lean{free\_threshold}, \lean{trueValue\_eq\_flow}, \lean{free\_threshold\_flow}
in \lean{LeanPaperCFreeBandV1}}

Stated for laws alone, the threshold is close to a definition. The flow form is what makes it a
computation: one parametric program over at most $\prod_\tau(n_\tau+1)$ states per layer. On
$g_m$ Theorems~\ref{thm:even} and~\ref{thm:odd} give it in closed form:
$\eta_{\rm free}=\frac{m-2}{2(m-1)}$ in even size and $1$ in odd size.

\paragraph{The additive price} Let $\Delta_g(\eta)=V_g(\eta)-V^{\rm free}_g$.

\begin{proposition}[The additive price is stable]
\label{prop:gap}
Let $\|g\|_1>0$ and $\eta\ge0$. Then $\Delta_g(\eta)\ge0$, $\Delta_g$ is nonincreasing, and
$\Delta_g(\eta)=0$ exactly when $\eta\ge\eta_{\rm free}(g)$. For all $g,g'$ with
$\|g\|_1,\|g'\|_1\ge\nu>0$,
\[
|\Delta_g(\eta)-\Delta_{g'}(\eta)|\le6\|g-g'\|_1/\nu.
\]
\end{proposition}
\leanbox{\lean{gap\_nonneg}, \lean{gap\_antitone}, \lean{gap\_zero\_iff}, \lean{gap\_stable} in
\lean{LeanPaperCFreeBandV1}}

The ratio $V_g(0)/V^{\rm free}_g$ needs $V^{\rm free}_g>0$, and it can blow up when the free
value is small. The additive price needs no such condition. In the exploration below, $96$ of the
$240$ count vectors have $V^{\rm free}_g=0$ to solver precision \NUM. There the ratio is undefined and the
additive price is not.

\paragraph{One item of each type, at every ratio}
\begin{theorem}[Two magnitudes, one item each]
\label{thm:twomag}
Let $g=(a,-a,b,-b)$ with $0<b<a$, so $D=\|g\|_1=2(a+b)$. For every $\eta\ge0$,
\[
V_g(\eta)=\frac{a-b}{2(a+b)}\,\max\Bigl\{\frac{1-\eta}2,\ \frac13\Bigr\}.
\]
So the frontier has one breakpoint, at $\eta=\frac13$, and $V_g(0)/V_g^{\rm free}=\frac32$ at
every ratio $a/b$.
\end{theorem}
\leanbox{\lean{base\_one}, \lean{base\_two}, \lean{base\_three}, \lean{per\_order},
\lean{lower\_band}, \lean{lower\_free}, \lean{law\_feasible}, \lean{law\_obj},
\lean{twoMag\_frontier}, \lean{twoMag\_price} in \lean{LeanPaperCTwoMagFourV1}}

Both sides are explicit.
\begin{itemize}
\item \emph{Dual.} The band at the first position puts an $a$-item first with probability at
least $(1-\eta)/2$, and that costs $a-b$ at budget one. On every order the three budgets
together cost at least $a-b$: either an $a$-item is first or last, or the first two items are one
$b$-item and one $a$-item.
\item \emph{Primal.} Put weight $\frac{1-\eta}2$ on the orders with the $a$-pair first,
$\frac{1-\eta}2$ on those with the $b$-pair first, and $\eta$ on
$(\pm b,\mp a,\pm a,\mp b)$, with two orientations each. For $\eta\le\frac13$ this law is
band-feasible with budget costs $\frac{a-b}{D}(\frac{1-\eta}2,\eta,\frac{1-\eta}2)$.
\end{itemize}
At $a/b=3$ this is the case $m=4$ of Theorem~\ref{thm:even}.

\paragraph{The two-magnitude family, numerically}
For $g=(a^{n_{+a}},(-a)^{n_{-a}},b^{n_{+b}},(-b)^{n_{-b}})$ with $a>b>0$, we solved the flow
program of Theorem~\ref{thm:typedflow} in floating point. The grid had $a/b\in\{3/2,2,3,5\}$,
counts at most $2$ per type with $m\ge3$ ($240$ vectors), and $\eta$ in steps of $0.05$ on
$[0,2]$. The script checks itself: the program reproduces Theorems~\ref{thm:even}
and~\ref{thm:odd} on $g_m$ for $m=4,\dots,9$ (largest gap $2.4\cdot10^{-16}$), $\eta=m$ gives the
unbanded value, and every frontier is nonincreasing. Breakpoints are read from second
differences on the grid, so they are conjectures and not results. \NUM{}
\begin{itemize}
\item Counts $(1,1,1,1)$: one breakpoint between $0.30$ and $0.35$, and
$V(0)/V^{\rm free}=\frac32$ at every tested ratio. Theorem~\ref{thm:twomag} proves this at
every ratio.
\item Counts $(1,1,2,2)$: price $\frac43$ with the band free from $\eta=1$ at $a/b\le2$, and
price $\frac53$ with the band free from $\eta=0.4$ at $a/b\ge3$.
\item Counts $(2,2,1,1)$: price $\frac53$, free from $0.8$, at every ratio.
\item Counts $(2,2,2,2)$: price $1.125$ at $a/b\le2$ and $1.75$ at $a/b\ge3$, each with two
kinks one grid step apart.
\end{itemize}
A second run refines the ratio to steps of $0.05$ on $[2,3]$, with $\eta$ in steps of
$0.025$. It reproduces the stored rows at $a/b=2,3$ exactly. It shows no single transition.
\NUM{}
\begin{itemize}
\item Counts $(1,1,2,2)$: the price falls from $\frac43$ at $a/b=2$ to $\frac76$ at $a/b=\frac52$.
It then rises, linearly from about $a/b=2.6$, to $\frac53$ at $a/b=3$. The free threshold moves
from $1$ down to $0.25$ and back to $0.4$.
\item Counts $(2,2,2,2)$: the price rises from $\frac98$ to a local maximum near $a/b=2.2$,
dips, and then rises linearly to $\frac74$.
\end{itemize}
So the frontier depends on the magnitude ratio piecewise, with several breakpoints in $a/b$ as
well as in $\eta$. Parity of $m$ alone does not decide the free-band gap: counts $(1,1,2,2)$ and
$(2,2,1,1)$ both have $m=6$. Proving these region laws needs a primal law and a matching dual
certificate per region; that is open.

\section{Robustness, and where it ends}
\label{sec:robust}

The family $g_{m}$ is highly structured, and it is fair to ask whether the factor-$m$ price
depends on exact cancellation. Relaxing the constraint is Theorem~\ref{thm:odd}. Perturbing
the weights is this section. The two are not combined.

For arbitrary $g'\in\mathbb{R}^{m}$, let $V(g';\eta)$ be the physical optimum. It
is~\eqref{eq:value} over every order law with the per-item band~\eqref{eq:band}, with the
baseline $b_{k}(g')$ computed from its definition and $D=\|g'\|_{1}$. Perturbations that break
the equal-weight groups are in scope, and $V(g_{m};\eta)=V_{m}(\eta)$
by~\eqref{eq:baseline}. The band is vacuous at $\eta=m$, so $V(g';m)$ is the unconstrained
optimum of $g'$ and $V(g';0)/V(g';m)$ is its price ratio. The normaliser cancels in the
ratio.

\begin{theorem}[Envelope]
\label{thm:robust}
Let $m\ge3$, $g'\in\mathbb{R}^{m}$, $\varepsilon=\|g_{m}-g'\|_{1}$ and
$0<d\le\min(m+4,\|g'\|_{1})$. Then $|V(g';\eta)-V_{m}(\eta)|\le\frac{3\varepsilon}{d}$ for
every $\eta\ge0$. Let $m$ be odd. Whenever $V(g';m)>0$ and $\frac{3\varepsilon}{d}\le V_{m}(0)$,
\[
\frac{V(g';0)}{V(g';m)}\ \ge\ \frac{V_{m}(0)-3\varepsilon/d}{V^{\mathrm{free}}_{m}+3\varepsilon/d},
\]
and for $\varepsilon\le\frac{1}{2m}$, with no further condition,
$\frac{V(g';0)}{V(g';m)}\ge\frac{m-1}{4}$.
\end{theorem}

\leanbox{\lean{family\_value\_robust}, \lean{even\_frontier\_robust} in
\lean{LeanPaperCStabilityFamilyV1}, from \lean{trueValue\_stable} in
\lean{LeanPaperCStabilityEndgameV1}; \lean{physical\_price\_ratio}, \lean{physical\_price\_linear}
in \lean{LeanPaperCRatioJoinV1}}

The constant $3$ pays for three moving parts: the prefix sum, the recomputed baseline and the
normaliser. The envelope alone cannot say where the price collapses, because at fixed
$\varepsilon$ its floor is bounded in $m$. The next theorem combines a lower bound valid for every perturbation with an upper bound
achieved by one explicit perturbation. The distinction between those quantifiers matters.

\begin{theorem}[Collapse scale, odd $m$]
\label{thm:collapse}
Let $m\ge3$ be odd.
\begin{enumerate}
\item[(a)] \emph{Upper side.} Take $\delta\in(0,1]$ and move one light item from $-1$ to
$-1+\delta$. This perturbation has $\|g_{m}-g'\|_{1}=\delta$ and a positive unconstrained
value, and its price ratio is at most $\frac{10}{\delta}$.
\item[(b)] \emph{Lower side.} Every $g'$ with $\varepsilon=\|g_{m}-g'\|_{1}\in[\frac1m,\frac16]$
and $V(g';m)>0$ has price ratio at least $\frac{1}{16\varepsilon}$. Every $g'$ with
$\varepsilon\le\frac1m$ and $V(g';m)>0$ has price ratio at least $\frac{m-1}{8}$. So on the
whole range $0<\varepsilon\le\frac16$ the ratio is at least
$\frac1{16}\min\bigl(m-1,\frac1\varepsilon\bigr)$, with no gap in $\varepsilon$.
\end{enumerate}
\end{theorem}

\leanbox{\lean{collapse\_upper}, \lean{collapse\_lower} in \lean{LeanPaperCCollapseScaleV1};
\lean{window\_arith}, \lean{collapse\_window}, \lean{collapse\_all},
\lean{collapse\_window\_satisfiable} in \lean{LeanPaperCCollapseWindowV1}}

Define the price that is guaranteed throughout a ball by
\begin{equation}
\mathcal R_m(\varepsilon)=
\inf_{\substack{\|g'-g_m\|_1\le\varepsilon\\V(g';m)>0}}
\frac{V(g';0)}{V(g';m)}.
\label{eq:robustball}
\end{equation}
The positivity condition avoids division by a zero free optimum. The set is nonempty since
it contains $g_m$.

\begin{corollary}[Sharp order of the uniform robustness guarantee]
\label{cor:robustball}
For odd $m\ge3$ and $0<\varepsilon\le1/6$,
\begin{align}
\frac1{16}\min\{m-1,\varepsilon^{-1}\}
&\le \mathcal R_m(\varepsilon)\nonumber\\
&\le\min\left\{\frac{(m+3)(m-1)}{2m},\frac{10}{\varepsilon}\right\}\nonumber\\
&\le10\min\{m-1,\varepsilon^{-1}\}.
\label{eq:robustballrate}
\end{align}
\end{corollary}
\begin{proof}
At a distance $d\in(0,\varepsilon]$, Theorem~\ref{thm:collapse}(b) bounds the ratio below by
$\min\{m-1,1/d\}/16\ge\min\{m-1,1/\varepsilon\}/16$.
At distance zero the exact ratio is $(m+3)(m-1)/(2m)$, which also satisfies that bound.
This lower bound is uniform, so it passes to the infimum.
For the upper bound use, separately, $g_m$ itself and the one-item perturbation in
Theorem~\ref{thm:collapse}(a) with $\delta=\varepsilon$. Both are feasible in the ball and
have positive free value. Finally $(m+3)/(2m)\le1$ for $m\ge3$.
\end{proof}
\leanbox{\lean{robustPrice}, \lean{member\_lower}, \lean{robust\_lower}, \lean{robust\_upper},
\lean{robust\_ball}, and for the rescaling example \lean{trueValue\_smul},
\lean{rescaled\_keeps\_price}, in \lean{LeanPaperCRobustBallV1}; inputs
\lean{physical\_price\_linear} in \lean{LeanPaperCRatioJoinV1}, \lean{collapse\_upper} in
\lean{LeanPaperCCollapseScaleV1}, \lean{collapse\_all} in \lean{LeanPaperCCollapseWindowV1}.}

\paragraph{Why this is not a law for every perturbation}
For any $\varepsilon>0$, take $g'=(1+\varepsilon/(m+4))g_m$.
Then $\|g'-g_m\|_1=\varepsilon$, but every prefix magnitude, every budget baseline and the
normaliser scale by the same positive number. The feasible laws do not change, so
$V(g';\eta)=V_m(\eta)$ at every $\eta$. Its price stays proportional to $m$ even at a fixed
nonzero radius. Thus~\eqref{eq:robustballrate} identifies when the guarantee can deteriorate,
not when every instance must deteriorate. For $\varepsilon\le1/(2m)$ the stronger uniform
lower bound $(m-1)/4$ follows directly from Theorem~\ref{thm:robust}.

The upper side is a parity argument. Every prefix sum of $g_{m}$ has the parity of its
length, which is what makes the free value small. After the move, the best drop of odd size
costs at most $1-\delta$ (drop the moved item), and the best drop of even size still costs
$0$ (leave it out). An order then pays at least $\delta$ at every odd budget before the moved
item and at every even budget from it on. Each pair of budgets $(2i+1,2i+2)$ contains one of
the two. So every order pays at least $\delta\frac{m-1}{2}$ in total, and the free value is at
least $\frac{\delta}{2\|g'\|_{1}}$. The envelope bounds the zero-band value from above.

\paragraph{What a search finds}
We searched for the worst perturbation at $m=5,6,7$, with random restarts and coordinate
descent over the $m!$ formulation and with baseline and normaliser recomputed. A minimising search returns
attained ratios, hence upper bounds on the infimum over its perturbation domain, not lower
certificates for that infimum. At $m=7$ the search keeps $43.5\%$
of the price at $\varepsilon=0.8$ and $30.1\%$ at $\varepsilon=1.6$. At $m=5,6$ no family in
the grid halved it. \NUM{} These sizes lie below the range where the constants of
Theorem~\ref{thm:collapse} separate the two regimes. The search is reported as a consistency
check, not as evidence of the scale.

\section{Choosing the budget after the draw}
\label{sec:expost}

If the budget is chosen after the order is drawn, the regret is
$\mathbb{E}_{P}[\max_{k}r_{k}]$, which is linear in $P$.

\begin{proposition}[Free ex-post value]
\label{prop:expost}
For every $m\ge3$, every order of $g_{m}$ has an interior budget with excess at least $2$,
and some order has excess at most $2$ at every budget. Hence
$\inf_{P}\mathbb{E}_{P}[\max_{0<k<m}r_{k}]=\frac{2}{m+4}=(m-1)V^{\mathrm{free}}_{m}$. For
every law, the ex-ante objective is at most the ex-post one.
\end{proposition}

\leanbox{\lean{pathValue\_eq}, \lean{path\_over\_free}, \lean{exists\_defect},
\lean{worstBudget\_even}, \lean{worstBudget\_odd},
\lean{objective\_le\_pathObjective} in \lean{LeanPaperCExPostV1}}

\begin{theorem}[Banded ex-post value]
\label{thm:expostband}
Let $V^{\mathrm{post}}_{m}(\eta)$ be the infimum of $\mathbb{E}_{P}[\max_{0<k<m}r_{k}]$ over laws
satisfying~\eqref{eq:band}. For every $\eta\ge0$,
\[
V^{\mathrm{post}}_{m}(\eta)=
\begin{cases}
\dfrac{2}{m+4}, & m\ge4\text{ even},\\[8pt]
\dfrac{2+2(1-\eta)_{+}/m}{m+4}, & m\ge3\text{ odd}.
\end{cases}
\]
\end{theorem}

\leanbox{\lean{expost\_even}, \lean{expost\_odd} in \lean{LeanPaperCExPostBandV1}}

Ex ante, uniformity costs a factor $\Theta(m)$ in odd size. Ex post, the free value is
already $m-1$ times larger, and the band adds at most a factor $1+\frac1m$. So the band price
and the timing price do not multiply.

The lower side is Proposition~\ref{prop:expost} plus one observation. An odd order with the
$+3$ last pays $4$ at budget $m-1$, and the band puts the $+3$ last with probability at least
$\frac{1-\eta}{m}$.

The upper side uses rotations. A uniform law on the $m$ cyclic shifts of one order has exactly
uniform marginals.
\begin{itemize}
\item Every shift of the flat even order $(-3,+1,-1,\dots,+1,-1,+3)$ pays at most $2$ at its
worst budget.
\item Every shift of the odd base order pays at most $2$, except the one that ends with $+3$,
which pays at most $4$.
\item The insertion mixture pays $2$ on every order at $\eta=1$, and blending covers $[0,1]$.
\end{itemize}

\section{Knowing the budget, and editing after it is revealed}
\label{sec:aware}

One order serves every budget at once. Two relaxations separate what that commitment costs,
and both keep the band. A policy that learns $k$ before acting may choose a different law at
each budget. A policy that commits to an order may still edit its prefix once $k$ is revealed.

\paragraph{The budget-aware value} A subset law $Q_{k}$ on $k$-sets meets the band when
$|\Pr_{Q_{k}}(i\in S)-k/m|\le\eta/m$ for every item. With
$r_{k}(S;g)=(|\sum_{i\in S}g_{i}|-b_{k}(g))/D$, let
\[
B_{g,k}(\eta)=\inf_{Q_{k}}\mathbb{E}_{Q_{k}}r_{k}(S;g),\qquad
B_{g}(\eta)=\max_{0<k<m}B_{g,k}(\eta),
\]
and let $A_{g}(\eta)=V_{g}(\eta)-B_{g}(\eta)$ be the value of knowing the budget.

\paragraph{Recourse} A plan commits to a law $P$ satisfying~\eqref{eq:band}. When $k$ is
revealed it replaces at most $R$ items of the prefix $S_{k}(\pi)$. Formally it is a joint law
of $(\pi,T)$ whose first marginal is $P$, with $|T|=k$, $|S_{k}(\pi)\setminus T|\le R$, and the
item marginals of $T$ in the band. Let $V^{(R)}_{g}(\eta)$ be the best worst-budget expected
regret of the final set $T$.

\begin{proposition}[Knowing, editing, committing]
\label{prop:sandwich}
For every $g$ with $\|g\|_1>0$, every $\eta\ge0$ and every $R\le R'$,
\[
0\le B_{g}(\eta)\le V^{(R')}_{g}(\eta)\le V^{(R)}_{g}(\eta)\le V^{(0)}_{g}(\eta)=V_{g}(\eta).
\]
If $\min(k,m-k)\le R$ at every interior budget, then $V^{(R)}_{g}=B_{g}$. In particular
$\lfloor m/2\rfloor$ exchanges always reach $B_{g}$, and one exchange reaches the optimum of the
budgets $1$ and $m-1$.
\end{proposition}
\leanbox{\lean{push\_feasible}, \lean{aware\_le\_true}, \lean{aware\_le\_recourse},
\lean{recourse\_mono}, \lean{recourse\_le\_true}, \lean{recourse\_zero}, \lean{sdiff\_card\_le},
\lean{recourse\_full} in \lean{LeanPaperCBudgetAwareV1}}

The prefix of a banded order is a banded subset law with the same regret, which gives
$B_g\le V^{(R)}_g$. Two $k$-sets differ by at most $\min(k,m-k)$ items, so a final set drawn
independently from a near-optimal $Q_{k}$ is always reachable. Lower bounds on a layer come
from linear certificates. If $\alpha+\sum_{i\in S}\beta_{i}\le D\,r_{k}(S;g)$ for every $k$-set
$S$, then
\[
B_{g,k}(\eta)\ge\Bigl(\alpha+\frac km\sum_{i}\beta_{i}-\frac\eta m\sum_{i}|\beta_{i}|\Bigr)\Big/D.
\]
\leanbox{\lean{linear\_lower}, \lean{layer\_ge\_of\_linear} in \lean{LeanPaperCBudgetAwareV1}}

\begin{theorem}[Knowing the budget is worth nothing on a band range]
\label{thm:zerovalue}
On $g_{m}$, $A_{g_m}(\eta)=0$ in two ranges:
\begin{itemize}
\item even $m\ge4$ and $0\le\eta\le\frac{m-2}{2(m-1)}$, where
$B_{g_m}(\eta)=V_{m}(\eta)=\frac{4(1-\eta)}{m(m+4)}$;
\item odd $m\ge3$ and $0\le\eta\le1-\frac{2}{m(m-1)}$, where
$B_{g_m}(\eta)=V_{m}(\eta)=\frac{m+3-(m+1)\eta}{m(m+4)}$.
\end{itemize}
\end{theorem}
\leanbox{\lean{aware\_even}, \lean{aware\_odd} in \lean{LeanPaperCBudgetGmV1}}

In each range a single budget carries the whole universal value. In even size, removing a heavy
item at budget one costs $2/D$ above the baseline, and the band gives each heavy item
probability at least $(1-\eta)/m$. In odd size, the dropped set at budget $m-1$ misses one item
$j$ and costs $|-1-g_{j}|/D$. The certificate there is $\beta=2$ on the light negatives, $-2$ on
the $+3$, and $0$ elsewhere. The two ranges end at the even breakpoint $\eta_*$ and at the odd
crossing point $\eta_t$ of Theorems~\ref{thm:even} and~\ref{thm:odd}.

\begin{theorem}[Four items: the value of the budget, and one exchange]
\label{thm:fourrecourse}
Let $g=(a,-a,b,-b)$ with $0<b<a$, and $c=\frac{a-b}{2(a+b)}$. For every $\eta\ge0$,
\begin{gather*}
B_{g}(\eta)=V^{(1)}_{g}(\eta)=c\max\Bigl\{\frac{1-\eta}2,0\Bigr\},\\
V^{(0)}_{g}(\eta)=V_{g}(\eta)=c\max\Bigl\{\frac{1-\eta}2,\frac13\Bigr\}.
\end{gather*}
So $A_{g}(\eta)=0$ on $[0,\frac13]$, $A_g(\eta)=\frac c2(\eta-\frac13)$ on $[\frac13,1]$, and
$A_g(\eta)=\frac c3$ for $\eta\ge1$. An exchange is needed exactly when $\eta>\frac13$, and one
exchange is enough.
\end{theorem}
\leanbox{\lean{aware\_twoMag}, \lean{advantage\_zero}, \lean{advantage\_mid},
\lean{advantage\_top}, \lean{recourse\_one\_twoMag}, \lean{recourse\_zero\_twoMag},
\lean{exchange\_needed} in \lean{LeanPaperCBudgetFamiliesV1}; \lean{twoMag\_frontier} in
\lean{LeanPaperCTwoMagFourV1}}

At budget $1$ the band forces probability at least $\frac{1-\eta}2$ that a heavy item is dropped,
and at budget $3$ that a heavy item is the one kept; each costs $c$. Giving each heavy item
probability $\frac{1-\eta}4$ and each light one $\frac{1+\eta}4$ meets the bound for $\eta\le1$. At budget two
the two cancelling pairs, each with probability $\frac12$, cost nothing even at $\eta=0$. The
plan commits to the uniform law on the four cyclic orders, which has exact uniform marginals.
At budgets $1$ and $3$ any two sets differ by one item. At budget two it keeps a cancelling
prefix and sends a mixed one, by one exchange, to a uniformly chosen cancelling pair.

\paragraph{Beyond the proved ranges, numerically} \NUM{} We solved the budget-aware layers, and
recourse through a transport between count vectors, by linear programming on $g_m$ for
$m=4,\dots,10$ and on two-magnitude multisets with counts $(1,1,2,2)$ and $(2,2,2,2)$, on the
grid $\eta\in\{0,0.01,\dots,1.2\}$ (\texttt{paperC\_budget\_recourse.json}, $6/6$ checks, among them the theorems above).
\begin{itemize}
\item On $g_m$, $A$ becomes positive at the first grid point past each proved range, so both
ranges appear sharp.
\item On every instance and grid point, one exchange reaches $B$.
\item With counts $(1,1,2,2)$ or $(2,2,2,2)$ and $a/b\le2$, knowing the budget helps already at
$\eta=0$.
\end{itemize}
The recourse values are for the count model. Lifting a count transport to exchanges of items,
by uniform relabelling within types, is not formalized.

\section{Evidence}
\label{sec:evidence}

Every number is produced by a script that checks itself against a statement proved in Lean,
and refuses to write output if a check fails. Figure~\ref{fig:ctheorems} draws the three
closed forms from the theorems. Its markers are floating-point LP solutions at $m\le9$; they
land on the formulas to $10^{-16}$ and certify nothing. Figure~\ref{fig:cstab} checks the
envelope against every stored perturbation.

\begin{table}[t]
\caption{Evidence packages. Each script self-checks against a Lean statement.}
\label{tab:evidence}
\centering
\footnotesize
\resizebox{\columnwidth}{!}{%
\begin{tabular}{@{}llr@{}}
\toprule
Package & Script & Checks \\
\midrule
Formula cross-check (ex ante, ex post) & \texttt{EXP\_paperC\_formula\_check} & 2/2 \\
Theorem figure & \texttt{FIG\_paperC\_theorems} & 3/3 \\
Flow program, breakpoints & \texttt{EXP\_paperC\_flow\_lp} & 5/5 \\
Breakpoint certificates & \texttt{EXP\_paperC\_breakpoint\_certs} & 5/5 \\
Exact Farkas roots & \texttt{EXP\_paperC\_exact\_dual} & 5/5 \\
Stored data vs.\ theorems & \texttt{FIG\_paperC\_endpoints} & 10/10 \\
Perturbation, both normalisations & \texttt{EXP\_paperC\_perturbation} & 9/9 \\
Adversarial price search & \texttt{EXP\_paperC\_robustness} & 4/4 \\
Window sums, value & \texttt{EXP\_paperC\_window} & 7/7 \\
\bottomrule
\end{tabular}}
\end{table}

\begin{figure*}[t]
\centering
\includegraphics[width=0.95\textwidth]{paperC_theorems.pdf}
\caption{Exact frontiers and robustness bounds, drawn from the theorems. Left: the whole frontier
$V_{m}(\eta)/V^{\mathrm{free}}_{m}$ (Theorems~\ref{thm:even}, \ref{thm:odd}), odd sizes solid
and even sizes dashed. Centre: the banded ex-post value over the free ex-post value
$\frac{2}{m+4}$ (Theorem~\ref{thm:expostband}); it is at most $1+\frac1m$, and exactly $1$ for
even $m$. Markers: floating-point LP solutions at $m\le9$ (\NUM), cross-checks only. Right:
the erosion of a uniform guarantee at $m=101$, theorem curves only. The lower curves hold
for every admissible perturbation; the dashed upper curve is attained as a bound by one
construction, not imposed on every perturbation (Corollary~\ref{cor:robustball}). Black: the envelope floor of
Theorem~\ref{thm:robust} with $d=m+4-\varepsilon$. Green: the linear guarantee $\frac{m-1}{4}$
for $\varepsilon\le\frac{1}{2m}$. Blue: the lower side $\frac{1}{16\varepsilon}$ on
$[\frac1m,\frac16]$. Dashed: the one-item construction $\frac{10}{\varepsilon}$. Dotted: the
unperturbed ratio. (\texttt{FIG\_paperC\_theorems}, 3/3.)}
\label{fig:ctheorems}
\end{figure*}

\begin{figure*}[t]
\centering
\includegraphics[width=0.92\textwidth]{paperC_stability.pdf}
\caption{The envelope against stored perturbations; nothing was re-solved. Left: all 220
perturbations of the physical objective, 160 group perturbations over the flow LP and 60
per-item perturbations over the $m!$ order LP. They are drawn against $3\|g_{m}-g'\|_{1}/d$
with $d=\min(\|g_{m}\|_{1},\|g'\|_{1})$. A point in the shaded region would flag a
mismatch requiring investigation of the model, solver and theorem; none enters it, and the
largest drift/envelope ratio is $0.32$. The producer version associated with these saved
results could skip solver failures, so the 220 recorded rows were checked against the full
prescribed grid; none was missing. Centre: the same drifts divided by the envelope. Objectives with the
normaliser or the baseline frozen are drawn only to show which object moves. Right: price
ratios attained by the adversarial search (markers), which are upper bounds on the searched infimum,
against the fixed-normaliser floor $\frac{1}{4\varepsilon+4/(m-1)}$ (lines, odd $m$; its
status is stated in the supplement).}
\label{fig:cstab}
\end{figure*}

\section{Limitations}
\label{sec:limits}

\begin{enumerate}
\item[1)] \textbf{One family for closed forms.} Every closed-form value is for $g_{m}$
of~\eqref{eq:family}. Other weights enter through Theorems~\ref{thm:robust},
\ref{thm:collapse}, \ref{thm:typedflow} and~\ref{thm:freeband}, which give an exact program
and threshold but no closed form. On two magnitudes the frontier is proved for one item of each
type (Theorem~\ref{thm:twomag}); for larger counts the map is floating point on a grid, and its
region laws are unproved.
\item[2)] \textbf{Robustness is a guarantee over a ball.} Corollary~\ref{cor:robustball}
determines the smallest attainable ratio only up to absolute constants and excludes zero
free values. It does not characterise every perturbation, nor give an exact robustness curve.
The exact frontiers at $g_m$ and this perturbation law are different results.
\item[3)] \textbf{Value, not algorithm.} Theorem~\ref{thm:typedflow} gives a program with
$\prod_\tau(n_\tau+1)$ states per layer, polynomial only for a fixed number of types. A generic
perturbation makes every weight distinct, and then $d=m$.
\item[4)] \textbf{Two objectives, kept apart.} Every ex-ante statement concerns
$\max_{k}\mathbb{E}r_{k}$, and every ex-post one concerns $\mathbb{E}\max_{k}r_{k}$.
\item[5)] \textbf{Budget information and recourse.} The value of knowing the budget is exact on
$g_m$ only in the two ranges of Theorem~\ref{thm:zerovalue}, and on four items at every band.
Beyond them it is numerical, and so is the observation that one exchange suffices; the recourse
computation uses the count model. In $B_g$ the per-budget law is a law on subsets. That it can be
realised by a full order satisfying~\eqref{eq:band} at every budget is not formalized.
\end{enumerate}

\section{Conclusion}

For one distribution over orders serving all removal budgets, uniform item marginals have
an exact price that depends sharply on parity. The complete frontier and explicit attaining
laws show where this price disappears; the composition flow preserves the optimisation over
all orders rather than restricting it to a constructed menu. Two changes of quantifier have
different effects. Choosing the budget after the draw raises the free cost but makes the
marginal band almost free. Requiring a guarantee over perturbed weights preserves a linear
odd-size price only to radius of order $1/m$; the guaranteed price thereafter scales as
$1/\varepsilon$ for $\varepsilon\le1/6$, although individual perturbations can preserve the original ratio.
Knowing the budget in advance can be worth nothing: on the family, over a range of bands, one
budget alone forces the universal value. Once the band is loose enough, commitment has a price,
and on four items a single exchange after the budget is revealed recovers all of it.
These results separate marginal fairness, budget timing and robustness within one precisely
specified compression objective.

\section*{Companion material}
The supplement gives the full proof architecture of each theorem, every supporting Lean
declaration, and the frozen builds it was checked in. The claim map records, for each
headline, its model, its declaration, its scope and its evidence. Corrections made during
preparation are recorded in the project ledger, not in this text.

\begin{thebibliography}{00}

\bibitem{singh2018} A. Singh and T. Joachims, ``Fairness of exposure in rankings,'' in
\emph{Proc. 24th ACM SIGKDD Int. Conf. Knowledge Discovery and Data Mining}, 2018,
pp.~2219--2228.

\bibitem{celis2018} L. E. Celis, D. Straszak, and N. K. Vishnoi, ``Ranking with fairness
constraints,'' in \emph{Proc. 45th Int. Colloq. Automata, Languages, and Programming
(ICALP)}, 2018.

\bibitem{singh2021} A. Singh, D. Kempe, and T. Joachims, ``Fairness in ranking under
uncertainty,'' in \emph{Advances in Neural Information Processing Systems 34}, 2021,
pp.~11896--11908.

\bibitem{disser2023} Y. Disser, M. Klimm, K. Schewior, and D. Weckbecker, ``Incremental
maximization via continuization,'' in \emph{Proc. 50th Int. Colloq. Automata, Languages, and
Programming (ICALP)}, LIPIcs vol.~261, 2023, art.~47.

\bibitem{bienkowski2026} M. Bienkowski, J. Blikstad, J. Byrka, M. Costa, Y. Disser, and
A. Lutz, ``Incremental submodular maximization: Better than greedy,'' arXiv:2606.28558, 2026.

\bibitem{koike2022} T. Koike, L. Lin, and R. Wang, ``Joint mixability and notions of negative
dependence,'' arXiv:2204.11438, 2022.

\bibitem{freeman2020} R. Freeman, N. Shah, and R. Vaish, ``Best of both worlds: Ex-ante and
ex-post fairness in resource allocation,'' in \emph{Proc. 21st ACM Conf. Economics and
Computation (EC)}, 2020.

\bibitem{gupta2022} A. Gupta, V. Gurunathan, R. Krishnaswamy, A. Kumar, and S. Singla,
``Online discrepancy with recourse for vectors and graphs,'' in \emph{Proc. ACM-SIAM Symp.
Discrete Algorithms (SODA)}, 2022.

\bibitem{schrijver2003} A. Schrijver, \emph{Combinatorial Optimization: Polyhedra and
Efficiency}. Berlin, Germany: Springer, 2003.

\bibitem{mathlib2020} The mathlib Community, ``The Lean mathematical library,'' in
\emph{Proc. 9th ACM SIGPLAN Int. Conf. Certified Programs and Proofs (CPP)}, 2020,
pp.~367--381.

\bibitem{demoura2021} L. de Moura and S. Ullrich, ``The Lean 4 theorem prover and
programming language,'' in \emph{Proc. 28th Int. Conf. Automated Deduction (CADE)}, 2021,
pp.~625--635.

\end{thebibliography}

\end{document}
